\documentclass[10pt]{article}
\usepackage[margin=20mm]{geometry}
\usepackage{amsmath,amssymb,fontspec,unicode-math,xurl,hyperref}
\IfFontExistsTF{DejaVu Serif}{\setmainfont{DejaVu Serif}}{\setmainfont{Latin Modern Roman}}
\setmathfont{latinmodern-math.otf}
\setlength{\parindent}{0pt}\setlength{\parskip}{6pt}\sloppy
\begin{document}
Abstract: The proof of the central cone-universality result contains an invalid block-to-feature construction. Its exclusive coordinates have the same sign on both sides, so the claimed cross-block weight cancels rather than being reproduced. A finite example satisfies the cone constraints and gives the zero matrix in place of a nonzero target. The quantitative dimension and error-budget expressions also need correction. An explicit alternative construction recovers qualitative universality with a different bound; the defect is in a core derivation, not a proof that the model cannot be universal.\par
Author: AI-assisted technical review\par
Date:\par
Title: Technical Verification: PieClam: A Universal Graph Autoencoder Based on Overlapping Inclusive and Exclusive Communities\par
\# Paper identity and scope\par
PieClam: A Universal Graph Autoencoder Based on Overlapping Inclusive and Exclusive Communities. Daniel Zilberg; Ron Levie. ICML 2025, PMLR 267, I25-43754. The reviewed source is the frozen published proceedings PDF. Page numbers below are one-based physical PDF pages.\par
Sections 2–3 and Appendix A.1–A.4 were examined in detail, with the main empirical discussion and selected Appendix D material read to bound the consequences. Original pages 5–7 and 13–16 were visually inspected. The reviewed code snapshot is `948cf7f87b572d0cd09886648074cd7198f9d7ce`, https://github.com/danizil/PieClam. Only selected likelihood, message-passing, probability and distance routines were read statically. No author code was executed.\par
\# Definitions and the affected claim\par
Write \(f=(t,s)\), \(L=\operatorname{diag}(I,-I)\), and \[T^C=\{(t,s): |s_c|\le t_c\text{ for all }c\}.\] The decoder \(p_{nm}=1-\exp(-f_n^\top Lf_m)\) is a probability on this cone because \(t_{nc}t_{mc}-s_{nc}s_{mc}\ge0\) coordinate by coordinate. Theorem 3.8 claims that every target graph can be approximated in the regularized log-cut distance by a cone-constrained code whose dimension depends only on the error tolerance.\par
The proof first approximates \(\widetilde A=-\log(1-(1-\delta)A)\) by a symmetric nonnegative block matrix \(B\), with block values \(b_{ij}\). The essential next assertion (p. 16) is that its prescribed features satisfy \(f_n^\top Lf_m=b_{k_nk_m}\). Without this equality the weak-regularity approximation does not transfer to the probability decoder.\par
\# Problem 1: the block representation cancels the required edges\par
For cross-class channels, Appendix A.4 assigns \(t=\sqrt{b_{ij}/4}\) and \(s=-\sqrt{b_{ij}/4}\) at both channel indices \((K-1)i+j\) and \((K-1)j+i\). Consider \(K=2\) singleton classes and \[B=\begin{pmatrix}0&1\\1&0\end{pmatrix}.\] This is a nonnegative symmetric block matrix as required by Definition A.4. The diagonal channel values vanish. Both cross-channel indices equal \(3\), and the displayed rule gives both nodes the same nonzero pair \((t_3,s_3)=(1/2,-1/2)\). Both features lie in \(T^{K^2}\). Nevertheless, \[f_1^\top Lf_2=\frac14-\frac14=0\ne1=b_{12}.\] All decoded probabilities are zero, whereas the claimed representation would have cross probability \(1-e^{-1}\). The cut-norm error in the block matrix is \(2/2^2=1/2\), despite starting from an exact block target.\par
This is not restricted to a particular value of the weak-regularity class-count bound. For any \(K\ge2\), a zero-diagonal target with nonzero off-diagonal blocks makes every assigned feature satisfy \(s_n=-t_n\), and hence makes the entire Lorentz Gram matrix zero. The independent \(K=4\) calculation with off-diagonal entries one gives cut error \(3/4\). Thus the failure remains for larger permitted class counts. The index map also collides, but correcting only the indexing does not correct the same-sign cancellation.\par
The original page image confirms the two negative signs; the adjacent bipartite example on p. 4, which uses opposite exclusive signs, provides clear counterevidence to interpreting this as a necessary limitation of IeClam. The error is in the printed construction, not in the cone’s expressive capacity.\par
The dimension and tolerance bookkeeping cannot be used literally either. Theorems 3.7–3.8 print a minus sign before \(\log(\epsilon/2)^2\). At \(\epsilon=1/2\), these expressions give \(K=-69.1852\) and \(C=2^{-28}\), neither a code dimension. A charitable reading is to square \(-\log(\epsilon/2)\) instead. Even then, A.4 requests only \(\|\widetilde A-B\|_\square\le\epsilon\) and subsequently adds \(\epsilon/2\), which yields at best \(3\epsilon/2\), not the stated \(\epsilon\). The strict bound also needs slack. These are not used as isolated typographical complaints: the displayed expressions do not establish the central quantitative guarantee.\par
This is Major under the claim-level standard. Theorem 3.8 is the proof supporting the paper’s probability-valued universality contribution. A repairable invalid construction remains an invalid core derivation. The counterexample does not disprove the existence of another encoding, and the benchmark results do not depend on implementing this proof construction.\par
\# An explicit repair with a changed quantitative bound\par
Every nonnegative symmetric \(K\)-block matrix has a valid cone representation using \(C=K(K+1)/2\) inclusive/exclusive coordinate pairs. Give each diagonal block \(i\) one coordinate, with \(t=\sqrt{b_{ii}}\), \(s=0\) on class \(i\) and zero elsewhere. Give every unordered pair \(i<j\) a distinct coordinate, writing \(a_{ij}=\sqrt{b_{ij}/2}\). On class \(i\) put \((t,s)=(a_{ij},a_{ij})\), on class \(j\) put \((t,s)=(a_{ij},-a_{ij})\), and put zero elsewhere.\par
Every feature belongs to \(T^C\). A pair coordinate contributes zero within either class, \(2a_{ij}^2=b_{ij}\) between its two classes, and zero otherwise. A diagonal coordinate contributes exactly \(b_{ii}\) within class \(i\). Summing proves \(f_n^\top Lf_m=B_{nm}\) for every pair, including self-pairs. The finite checks for \(K=2,4,6\) agree within \(4.45\cdot10^{-16}\); the preceding argument is the general verification.\par
For precision, a repaired qualitative statement is: for every \(0<\epsilon<1\) and every symmetric binary adjacency matrix \(A\), there is a code in \(T^C\) with \(D_\square(P\|A)<\epsilon\) and a dimension bound independent of the number of nodes. One sufficient bound follows by setting \[\delta=\eta=\epsilon/3,\quad R=\log(3/\epsilon),\quad
 K\le 4^{\lceil R^2/\eta^2\rceil},\quad C=K(K+1)/2.\] Use the nonnegative block approximation to \(\widetilde A=-\log(1-(1-\delta)A)\) with cut error at most \(\eta\), then the exact encoding above. The decoded \(P=1-e^{-B}\) has entries in \([0,1)\) and \[D_\square(P\|A)\le\delta+\|\widetilde A-B\|_\square
 \le 2\epsilon/3<\epsilon.\] Smaller actual partitions can be padded by unused code coordinates. Tolerances at least one are covered by applying the construction at a smaller tolerance. This restores qualitative universality, with an explicitly changed dimension bound; it does not validate the printed bound or claim author adoption.\par
\# Other checks and boundaries\par
The unrestricted construction in A.3 represents a signed matrix as a difference of two Gram matrices. That algebra is valid. A negative entry can, however, give \(1-e^{-c}<0\), which is outside the probability-valued decoder definition. Section 3.5 explicitly discloses this possibility. Thus the first theorem should be read as a signed approximation result; it cannot replace the cone argument when a probabilistic decoder is required.\par
For Claim 3.6, averaging the four bipartite blocks and using \(|a^2-\langle q,y\rangle|\ge a^2-\langle q,y\rangle\) supports the claimed lower bound under its self-pair convention. The no-self-loop setting is distinguished in Appendix A.2; the bound is not silently transferred to a different convention here.\par
An independent finite-difference calculation verifies equation (11)’s likelihood gradient at an interior cone point, with maximum discrepancy \(1.62\cdot10^{-10}\) at step \(10^{-5}\). The inspected probability routine uses \(1-\exp(-f_n^\top Lf_m)\) and explicitly zeros its diagonal. The inspected log-cut routine estimates a cut norm at a supplied regularizer using rounded/SDP outputs. Figure 1 itself labels a fixed \(d=0.01\); it does not prove convergence in the infimum-based distance or global likelihood optimization. Main Figures 2–3 visually illustrate prior and block reconstruction, without independently verifying data provenance or accuracy.\par
\# Conclusion\par
Dimension | Recorded state\par
Highest supported impact | Major\par
Scientific adjudication | Pending\par
Repair verification | Qualitative theorem repaired; bound changed\par
Author adoption | Not established\par
Human reviews | None recorded\par
The highest supported impact is Major because the central cone-universality derivation is invalid as printed. A verified alternative construction restores the qualitative result with a changed dimension bound. No empirical failure or impossibility of the proposed decoder is established.\par
1 Daniel Zilberg; Ron Levie. PieClam: A Universal Graph Autoencoder Based on Overlapping Inclusive and Exclusive Communities. PMLR 267 (2025), published PDF. https://proceedings.mlr.press/v267/zilberg25a.html.\par
\# Reproducibility and limits\par
`independent\_check.py` implements only the printed and repaired finite constructions, dimension arithmetic and an independent likelihood-gradient diagnostic. It imports no author modules. Source reads included `clamiter.py` lines 142–219 and 1539–1604, and `utils/utils.py` lines 50–109 and 407–438. Their presence does not establish the provenance of historical runs. Other source branches, prior optimization, raw datasets, downstream results and external references were not exhaustively audited.\par
For completeness, the block approximation used in the repair can be justified by the usual finite energy argument. Normalize the matrix inner product by \(N^2\). Let \(B\) be the blockwise average of a matrix in \([0,R]^{N\times N}\) on a common row/column partition. If its cut error exceeds \(\eta\), refine every part by a witnessing pair of subsets \(U,V\), multiplying the number of classes by at most four. The refined projection \(B'\) satisfies \(\langle B'-B,1_U1_V^\top\rangle
=\langle A-B,1_U1_V^\top\rangle\). Since the normalized norm of the indicator is at most one, \(\|B'-B\|_2>\eta\). Orthogonality of the projections increases \(\|B\|_2^2\) by more than \(\eta^2\). This energy is at most \(R^2\), so fewer than \(R^2/\eta^2\) unsuccessful refinement steps are possible, giving the stated conservative class bound. Block averages remain nonnegative and symmetric. This argument verifies the repaired qualitative result without assuming that the published coordinate formula is correct.\par
This initial review was performed by the assigned AI review unit. No human review or independent acceptance is claimed. The matching verification directory contains physical-page extracts, source checks, independent calculations, compilation records, visual inspection and provenance. Source SHA-256: `07302eb2314bcea3792b6a3108df7a3190952dde90973adc68ec4f290a4e22fb`.\par
\end{document}
